\section{Output dataset requirements and implications thereof}
\label{sec:output-dataset-requirements}

The output dataset requirements are very similar to the requirements that apply to the historical dataset [TODO ref historical manuscript].

As for the historical dataset, the first requirement is to work within CMIP's overall delivery timeline.
The timeline for future forcing delivery, like historical forcing delivery, is relatively strict by scientific standards.
As a result, we weigh the benefits of improvements against delays that would result from making improvements in all cases
and some identified bugs are not fixed.

As for the historical dataset, the second requirement is to meet CMIP's overall experiment design.
For the future forcing experiments, this currently means that we need to produce concentrations in line with the emissions produced by the experiment designer.
To date, we have only done this for ScenarioMIP [cite van Vuuren] and PolMIP's VL-CF scenario [cite Annika?].
If we were to produce other future forcing datasets,
they would either be produced on the same specified emissions-basis
or we would need to work with the experiment provider to design a new method.
Supporting ScenarioMIP's needs drives the temporal extent of our datasets.
ScenarioMIP's experiments start in 2022 (immediately following the end of the historical experiment)
and their extensions run until 2500-12.
% value-check: {"tag": "output-first-year", "unit": "yr", "lower": 2022, "upper": 2022}
% value-check: {"tag": "output-last-year", "unit": "yr", "lower": 2500, "upper": 2500}
The dataset we present here meets these requirements, Sect.~\ref{ssec:prep-for-esgf}.
In theory, other timespans could be produced if experiment designs required this
(as above, we would have to work with experiment designers to support this if custom methods are needed).
As for the historical dataset, there should only be one forcing dataset
so forcing dataset producers must produce a single, best-estimate dataset,
rather than being able to produce a number of versions that span the uncertainty in their forcing dataset.

As for the historical dataset, the third requirement is that the data must support use as an input to ESM simulations.
In addition, it must support use as an input to ESM simulations that branch from a historical simulation.
As a result, the data must be complete i.e. not contain any missing values
and it must be `harmonised' i.e. branch smoothly from the historical dataset. [TODO consider switching from historical dataset to dataset used for the historical experiment throughout as it is more explicit]
The use within ESMs also has implications for the required accuracy of the data.
As for the historical simulation, 0.01~W/m^2 is a tight, but justifiable threshold for identifying which methods need to be considered more carefully.
This threshold should not be confused with the threshold for being able to distinguish between two future experiments,
i.e. for a reasonable radiative forcing spacing between experiments, that is a different thing.

Taken together, these considerations define the output requirements for this study:
a temporally (at least from 2022 to end of 2500),
spatially (global coverage) complete dataset
that is appropriate for ESM simulations
and consistent with the underlying future experiment design
(in this paper this means consistent with the underlying future experiment emissions).

The key implications of these requirements are as follows.
Firstly, like for the historical dataset, the requirement to produce a single forcing estimate forces us to pick a single future projection.
% we cannot produce a range.
For future experiments, this will almost inevitably mean that the greenhouse gas forcing dataset we provide
will differ from projections of greenhouse gas concentrations that are made based by one or more emissions-driven earth system models (ESMs).
This is entirely expected.
% If we awnt to labour the point
% we have to produce forcings in advance of ESM runs
% and we cannot use an ESM to produce the forcings, because that would be too slow.
% As a result, the forcings we produce will, almost by definition, differ from those that are made by ESMs.

Secondly, these forcing datasets are provided within the context of the wider experiment design.
Hence, all the limitations that apply to that wider experiemnt design,
for example the need to pick an illustrative set of experiments
rather than being able to pick a set which covers thousands of experiments,
also apply to this dataset.
This is not an exhaustive set of possible future greenhouse gas concentrations.
It is simply greenhouse gas concentrations in line with the experiments that are being run in CMIP7.
% We are, essentially, only translators of the underlying experiment design in this regard.

Finally, the requirement of harmonisation means that we need to use a method
which can be smoothly and justifiably extended from the past to the future.
In our case, this means that we need an underlying method that can justifiably be used to produce
greenhouse gas concentration forcing for both year 1 and year 2500.
This is a large driver behind our choice of method
and there aren't easy answers (Sects.~\ref{sec:methods} and \ref{sec:discussion}).

Given these requirements and their implications,
potential users should carefully consider whether this dataset is appropriate for their use case.
% Maybe mention other limits here, let's see
